In this section, I present some applications that can immediately benefit from this framework outside of software systems.

\subsection{Universal Design Language}

Designs are the universal structures used across many disciplines. The ability to unambiguously establish a closed semantic world and reason about it has applications well beyond software systems. If implemented with the right strategy, this framework has the ability to become a universal design language used in many disciplines.

Since the information contained in the DAL has unambiguous meaning established by the semantic toolkit, it can be used to establish a powerful visualization framework. This would allow engineers to visualize the closed semantic world in a meaningful way and gain insights into their design by reasoning about its meaning.

The visualization framework can also be transformed into an interactive framework, establishing the ability to interact with the closed semantic world and this can be used in many different applications.

\subsection{Digital Twins}

Since this framework extends the design into a computable semantic model (CSM), it extends the CSM to become a digital twin of what it is modelling.

In software systems, the design establishes a closed semantic world that mirrors what participates in reality because the design is what defines it (this statement is even more valid with CSM's constructed using CSP's). This is a unique property of software systems since they operate entirely on a stack that is symbolically defined by humans. However, in other disciplines, reality is modelled by engineers to be able to represent it accurately and predict it. In these applications, the CSM becomes the framework through which that modelling happens. 

In the digital twin models, information collected by sensors is used to inform the model and refine it so that it is a more accurate representation of reality and to improve its ability to make more accurate predictions. In these cases, the domain structure of the collected data is unambiguously established by the CSM, as such, the data can be domain compressed using CLP's engine. This extends the use of CSM and CLP beyond software systems. 

\subsection{Modelling Knowledge}

The applications of CSM are only limited by what can be modelled, anything that can be modelled as a closed semantic world can be reasoned about using the engine. 

A personal favorite of mine is to model concepts using the CSM so that the model can be used as a teaching tool. This applications exploits the fact that when humans teach, they are ultimately building a closed semantic world to convey the concept to another human mind. Using the CSM, the semantic world necessary to represent the idea can be transformed into a computable model in which every step in the process is informed and validated by the model. 

There are many interesting ways in which this can be used to turn teaching and learning into an unambiguous process that converges on complete understanding.

\subsection{DCS for Archiving using Custom Specification}

Since CLP will be transformed into a domain specific compression engine using this framework, it can also be extended into a compression tool for optimally compressing any dataset. This can be enabled by providing the compression specification for the dataset.

The DCS engine will be able to use the compression specification to search the compressed data without decompression. It will turn CLP into an efficient compression tool that can be configured using a specification and it can be used across many industries.